\begin{frame}{Проблема отказоустойчивости в RTOS}
\begin{itemize}
    \item В системах реального времени (RTOS) задачи могут завершаться с ошибками или переставать отвечать;
    \item Традиционные подходы могут приводить к остановке всей системы или влиять на другие задачи;
    \item Перезапуск задачи является простым и эффективным способом восстановления без влияния на другие задачи;
\end{itemize}
\end{frame}

\begin{frame}[fragile]{Механизм контроля и перезапуска задач}
\begin{verbatim}
void supervisor() {
    for (int i = 0; i < TASK_NUM; i++) {
        if (tasks[i].status == FAILED) {
            restart_task(tasks[i].id);
        }
    }
}
\end{verbatim}

\vspace{0.5em}
\textbf{Описание:}
\begin{itemize}
    \item Периодическая проверка состояния задач;
    \item Обнаружение сбоя;
    \item Перезапуск только одной задачи;
    \item Остальная система продолжает работу.
\end{itemize}

\vspace{0.5em}
\textbf{Реализация механизма мониторинга} может быть различной:  
\textit{heartbeat, error hook, link}.
\end{frame}

\begin{frame}[fragile]{Метод 1: Heartbeat}
\begin{verbatim}
while (1) {
    if (recv_heartbeat(pid) == OK)
        update_timestamp(pid);
    if (now - last_seen(pid) > TIMEOUT) {
        pok_proc_create(&attr, &pid);
    }
}
\end{verbatim}

\vspace{0.5em}
\begin{itemize}
    \item Контроль активности задачи через heartbeat;
    \item Тайм-аут $\rightarrow$ задача считается отказавшей;
    \item Выполняется перезапуск.
\end{itemize}

\vspace{0.5em}
\textbf{Плюсы:} простота, не требует изменения ядра  
\textbf{Минусы:} задержка обнаружения
\end{frame}

\begin{frame}[fragile]{Метод 2: Error Hook}
\begin{verbatim}
pok_error_hook(my_hook);

void my_hook(pid) {
    send(supervisor, pid);
}
\end{verbatim}

\vspace{0.5em}
\begin{itemize}
    \item Срабатывает при ошибке задачи;
    \item Уведомляет механизм контроля;
    \item Выполняется восстановление.
\end{itemize}

\vspace{0.5em}
\textbf{Плюсы:} высокая скорость реакции  
\textbf{Минусы:} зависит от поддержки системы
\end{frame}

\begin{frame}[fragile]{Метод 3: Link-механизм (ядро)}
\begin{verbatim}
pok_sys_link(pid);
\end{verbatim}

\vspace{0.5em}
\begin{itemize}
    \item Ядро хранит связи между задачами;
    \item При завершении задачи отправляется сигнал;
    \item Выполняется автоматическое восстановление.
\end{itemize}

\vspace{0.5em}
\textbf{Плюсы:} максимальная скорость реакции  
\textbf{Минусы:} требуется модификация ядра
\end{frame}

\begin{frame}{Сравнение методов}
\begin{center}
\begin{tabular}{|c|c|c|c|}
\hline
\textbf{Метод} & \textbf{Сложность} & \textbf{Реакция} & \textbf{Изменение ядра} \\
\hline
Heartbeat & низкая & средняя (мс) & нет \\
\hline
Error Hook & средняя & высокая (мкс) & зависит \\
\hline
Link & высокая & мгновенная & да \\
\hline
\end{tabular}
\end{center}

\vspace{0.8em}
\end{frame}
